\chapter{Fundamental Theorems for Banach Spaces}

\section{Hahn--Banach theory}

\begin{notedefinition}[Linear functional]
A linear functional on a vector space $V$ is a linear map
$\varphi\colon V\to\mathbb K$.
\end{notedefinition}

\begin{notetheorem}[Bounded functionals and closed kernels]
Let $V$ be a normed space and let $0\ne\varphi\colon V\to\mathbb K$ be
linear.  The following are equivalent:
\begin{enumerate}
  \item $\varphi$ is bounded;
  \item $\varphi$ is continuous;
  \item $\Ker\varphi$ is closed.
\end{enumerate}
Since $\varphi\ne0$, its kernel is automatically a proper subspace.  That
properness alone does not imply boundedness.
\end{notetheorem}

\begin{proof}
Boundedness and continuity are equivalent, and continuity makes the inverse
image of $\{0\}$ closed.  Suppose $\Ker\varphi$ is closed.  Choose
$x_0\notin\Ker\varphi$.  Every $x$ has the unique form
$x=y+\alpha x_0$ with $y\in\Ker\varphi$, and closedness gives
$d:=\dist(x_0,\Ker\varphi)>0$.  Hence
\[
  d|\alpha|
  \le\|\alpha x_0-y'\|
  \le\|x\|
\]
for a suitable $y'\in\Ker\varphi$, while
$|\varphi(x)|=|\alpha|\,|\varphi(x_0)|$.  Thus
$|\varphi(x)|\le |\varphi(x_0)|d^{-1}\|x\|$.
\end{proof}
\sourcepages{46--47}

\begin{notetheorem}[Bases and discontinuous functionals]
Every vector space has a Hamel basis.  Consequently, every
infinite-dimensional normed space has a discontinuous linear functional.
\end{notetheorem}

\begin{proof}
The basis theorem follows from Zorn's lemma.  Choose linearly independent unit
vectors $(e_n)$ and extend them to a Hamel basis.  Prescribe
$\varphi(e_n)=n$ and $\varphi=0$ on the remaining basis vectors.  The
resulting algebraic linear functional is unbounded on the unit sphere.
\end{proof}
\sourcepages{47--48}

\begin{notelemma}[One-dimensional extension step]
Let $V$ be a real normed space, let $U\subseteq V$ be a subspace, let
$\psi\in U^*$, and let $h\notin U$.  There is an extension
$\widetilde\psi\colon U+\mathbb Rh\to\mathbb R$ with
$\|\widetilde\psi\|=\|\psi\|$.
\end{notelemma}

\begin{proof}
Any extension has the form
\[
  \widetilde\psi(u+th)=\psi(u)+tc.
\]
The estimate $|\widetilde\psi(u+th)|\le\|\psi\|\|u+th\|$ is equivalent to
choosing $c$ in the interval
\[
\sup_{u\in U}\bigl(-\|\psi\|\|u+h\|-\psi(u)\bigr)
\le c\le
\inf_{v\in U}\bigl(\|\psi\|\|v+h\|-\psi(v)\bigr).
\]
For $u,v\in U$, the left expression at $u$ is at most the right expression at
$v$ because
\[
  \psi(v)-\psi(u)=\psi(v-u)
  \le\|\psi\|\|v-u\|
  \le\|\psi\|\bigl(\|u+h\|+\|v+h\|\bigr).
\]
Thus the interval is nonempty.
\end{proof}
\sourcepage{49}

\begin{notetheorem}[Hahn--Banach theorem]
Let $X$ be a real vector space, let $Y\subseteq X$ be a subspace, and let
$p\colon X\to\mathbb R$ be sublinear.  If $f\colon Y\to\mathbb R$ is linear
and $f(y)\le p(y)$ on $Y$, then $f$ has a linear extension
$F\colon X\to\mathbb R$ satisfying $F(x)\le p(x)$ for all $x\in X$.
\end{notetheorem}

\begin{proof}
Partially order all dominated extensions by extension.  Unions of chains are
again dominated extensions, so Zorn's lemma gives a maximal one.  If its
domain $M$ is not all of $X$, choose $x_0\notin M$.  An extension to
$M+\mathbb Rx_0$ has the form
$F_\alpha(m+tx_0)=F(m)+t\alpha$.  Domination is equivalent to
\[
  \sup_{m\in M}\bigl(F(m)-p(m-x_0)\bigr)
  \le\alpha\le
  \inf_{m\in M}\bigl(p(m+x_0)-F(m)\bigr).
\]
Sublinearity shows the left endpoint does not exceed the right endpoint,
contradicting maximality.  Hence $M=X$.
\end{proof}
\sourcepages{50,53--54}

\begin{notetheorem}[Complex Hahn--Banach theorem]
Let $X$ be a complex vector space, let $Y\subseteq X$, and let
$p\colon X\to[0,\infty)$ be a seminorm.  If $f\colon Y\to\mathbb C$ is
complex linear and $|f(y)|\le p(y)$, then there is a complex-linear extension
$F\colon X\to\mathbb C$ with $|F(x)|\le p(x)$.
\end{notetheorem}

\begin{proof}
Let $u=\operatorname{Re}f$, regarded as a real-linear functional.  Extend
$u$ by the real theorem to a real-linear $U$ with $|U(x)|\le p(x)$.  Define
\[
  F(x)=U(x)-iU(ix).
\]
Then $F$ is complex linear and extends $f$.  Given $x$ with $F(x)\ne0$,
choose $\theta$ so that $e^{-i\theta}F(x)=|F(x)|$.  Since
$F(e^{-i\theta}x)=e^{-i\theta}F(x)$,
\[
  |F(x)|=\operatorname{Re}F(e^{-i\theta}x)=U(e^{-i\theta}x)
  \le p(e^{-i\theta}x)=p(x).
\]
\end{proof}
\sourcepages{50--51,55--56}

\begin{notecorollary}[Norm-preserving extension]
Let $X$ be a real or complex normed space, let $Y\subseteq X$, and let
$f\in Y^*$.  Then $f$ has an extension $F\in X^*$ satisfying
\[
  \|F\|=\|f\|.
\]
\end{notecorollary}

\begin{proof}
Apply Hahn--Banach with the sublinear functional
$p(x)=\|f\|\|x\|$.  The domination yields $\|F\|\le\|f\|$; restriction to
$Y$ yields the reverse inequality.
\end{proof}
\sourcepage{56}

\begin{notetheorem}[Norming functional]
For every nonzero $x_0$ in a normed space $X$, there is $f\in X^*$ such that
\[
  \|f\|=1,
  \qquad
  f(x_0)=\|x_0\|.
\]
Consequently,
\[
  \|x\|=\sup_{0\ne f\in X^*}\frac{|f(x)|}{\|f\|}
       =\sup_{\|f\|\le1}|f(x)|.
\]
\end{notetheorem}

\begin{proof}
On $\Span\{x_0\}$ define
$g(\alpha x_0)=\alpha\|x_0\|$.  It has norm one.  Extend it by
Hahn--Banach.  The norm formula follows from the general estimate
$|f(x)|\le\|f\|\|x\|$ and this norming functional.
\end{proof}
\sourcepage{57}

\begin{notetheorem}[Taylor--Foguel unique-extension criterion]
A normed space $X^*$ is strictly convex if and only if, for every subspace
$Y\subseteq X$ and every $y^*\in Y^*$, there is at most one norm-preserving
extension of $y^*$ to $X$.
\end{notetheorem}

\begin{proof}[Corrected proof outline]
If $F_1\ne F_2$ are norm-preserving extensions of $y^*$, then
$(F_1+F_2)/2$ is also an extension and has norm at least $\|y^*\|$.
Strict convexity gives a contradiction after normalization.

Conversely, suppose $X^*$ is not strictly convex.  Choose distinct
$F_1,F_2\in S_{X^*}$ with $\|(F_1+F_2)/2\|=1$, and put
$H=(F_1+F_2)/2$ and $Y=\Ker(F_1-F_2)$.  Choose $x_n\in S_X$ and unimodular
scalars so that, after replacing $x_n$ by those scalar multiples,
$H(x_n)\to1$.  Since $\operatorname{Re}F_j(x_n)\le1$ and their average
tends to $1$, both $F_1(x_n)$ and $F_2(x_n)$ tend to $1$.  Hence
$(F_1-F_2)(x_n)\to0$.  For the nonzero functional $F_1-F_2$,
\[
  \dist(x,Y)=\frac{|(F_1-F_2)(x)|}{\|F_1-F_2\|},
\]
so there are $y_n\in Y$ with $\|x_n-y_n\|\to0$.  It follows that
$\|H|_Y\|=1$.  But $F_1|_Y=F_2|_Y=H|_Y$, and $F_1,F_2$ are two distinct
norm-preserving extensions of this common restriction.
\end{proof}
\sourcepages{58--59}

\begin{notetheorem}[Separating a point from a closed subspace]
Let $Y$ be a closed subspace of a normed space $X$, let $x_0\notin Y$, and
put $d=\dist(x_0,Y)$.  Then there is $f\in X^*$ such that
\[
  \|f\|=1,
  \qquad
  f|_Y=0,
  \qquad
  f(x_0)=d.
\]
\end{notetheorem}

\begin{proof}
On $Y+\mathbb Kx_0$ define $g(y+\alpha x_0)=\alpha d$.  Since
$d\le\|x_0-y\|$ for every $y\in Y$,
\[
 |g(y+\alpha x_0)|=|\alpha|d
 \le\|y+\alpha x_0\|.
\]
Thus $\|g\|\le1$, while approximating the distance shows $\|g\|=1$.
Hahn--Banach extends $g$ with the same norm.
\end{proof}
\sourcepage{60}

\begin{notecorollary}[Annihilator criterion for density]
For a subspace $Y\subseteq X$,
\[
  \overline Y=X
  \quad\Longleftrightarrow\quad
  \{f\in X^*:f|_Y=0\}=\{0\}.
\]
\end{notecorollary}

\begin{proof}
If $Y$ is dense, continuity forces every functional vanishing on $Y$ to
vanish on $X$.  If $\overline Y\ne X$, separate a point outside
$\overline Y$ from the closed subspace $\overline Y$.
\end{proof}

\begin{notetheorem}[A separable dual forces a separable space]
If $X^*$ is separable, then $X$ is separable.
\end{notetheorem}

\begin{proof}
Choose a dense sequence $(f_n)$ in $S_{X^*}$ and, for each $n$, choose
$x_n\in S_X$ with $|f_n(x_n)|>1/2$.  Let
$M=\overline{\Span\{x_n:n\ge1\}}$.  If $M\ne X$, Hahn--Banach gives
$f\in S_{X^*}$ vanishing on $M$.  Choose $n$ with $\|f-f_n\|<1/2$.
Then $|f_n(x_n)|=|(f_n-f)(x_n)|<1/2$, a contradiction.
\end{proof}
\sourcepages{61--62}

\subsection{Geometric Hahn--Banach theorems}

\begin{notedefinition}[Minkowski functional]
Let $C\subseteq X$ be open, convex, and contain $0$.  Its gauge is
\[
  p_C(x)=\inf\{t>0:x\in tC\}.
\]
Then $p_C$ is sublinear, continuous, and
\[
  C=\{x:p_C(x)<1\}.
\]
\end{notedefinition}

\begin{proof}
Openness and absorption make $p_C$ finite.  Positive homogeneity follows by
rescaling.  If $x\in(r+\varepsilon)C$ and
$y\in(s+\varepsilon)C$, convexity gives
$x+y\in(r+s+2\varepsilon)C$, so
$p_C(x+y)\le p_C(x)+p_C(y)$.  A norm ball contained in $C$ gives
$p_C(x)\le M\|x\|$, hence continuity.  The displayed description of $C$
follows from openness and convexity.
\end{proof}

\begin{notetheorem}[Separation of two open convex sets]
Let $A,B$ be nonempty, disjoint, open, convex subsets of a real normed space
$X$.  Then there are a nonzero $f\in X^*$ and $\gamma\in\mathbb R$ such that
\[
  f(a)\le\gamma\le f(b)
  \qquad(a\in A,\ b\in B).
\]
\end{notetheorem}

\begin{proof}
Put $C=A-B$.  Then $C$ is open and convex and $0\notin C$.  Choose
$c_0\in C$, set $U=C-c_0$, and let $p_U$ be its Minkowski functional.  The
point $x_0=-c_0$ does not belong to $U$, so $p_U(x_0)\ge1$.  On
$\Span\{x_0\}$ define $g(tx_0)=t p_U(x_0)$ and extend it by Hahn--Banach
to a continuous functional $f\le p_U$.  For $c\in C$ one has
$c-c_0\in U$, and therefore
\[
  f(c)=f(c-c_0)+f(c_0)<1-p_U(x_0)\le0.
\]
Thus $f(a)<f(b)$ for every $a\in A$ and $b\in B$.  Consequently
$\sup_A f\le\inf_B f$; choose $\gamma$ between these two numbers.
\end{proof}
\sourcepages{63--65}

\begin{notetheorem}[Strict separation of a compact set]
Let $A$ be closed and convex and let $K$ be compact and convex in a real
normed space, with $A\cap K=\varnothing$.  Then there are
$f\in X^*\setminus\{0\}$, $\gamma\in\mathbb R$, and $\delta>0$ such that
\[
  f(a)\le\gamma-\delta<\gamma+\delta\le f(k)
  \qquad(a\in A,\ k\in K).
\]
\end{notetheorem}

\begin{proof}
Compactness and closedness give $d=\dist(A,K)>0$.  The open convex
neighbourhoods $A+B(0,d/3)$ and $K+B(0,d/3)$ are disjoint.  Separate them by
the preceding theorem.  The two added balls produce a positive gap in the
functional values.
\end{proof}
\sourcepage{65}

\begin{notetheorem}[Nonunique extensions over a strict closed subspace]
Let $Y$ be a proper closed subspace of a normed space $X$ and let
$\ell\in Y^*$.  Besides a norm-preserving extension, there are continuous
extensions with strictly larger norm.
\end{notetheorem}

\begin{proof}
Choose a norm-preserving extension $L\in X^*$.  Since $Y$ is proper and
closed, Hahn--Banach gives a nonzero $g\in X^*$ with $g|_Y=0$.  Then
$L+tg$ extends $\ell$ for every scalar $t$, and
$\|L+tg\|\to\infty$ as $|t|\to\infty$.
\end{proof}
\sourcepage{66}

\begin{noteexercise}[Continuity of sublinear functionals]
If a sublinear functional $p$ on a normed space is continuous at $0$ and
$p(0)=0$, then it is continuous everywhere.
\end{noteexercise}

\begin{proof}
Subadditivity gives
\[
  p(x)-p(y)\le p(x-y),
  \qquad
  p(y)-p(x)\le p(y-x).
\]
Both right-hand sides tend to zero as $x\to y$.
\end{proof}
\sourcepage{67}

\begin{noteexercise}[Kernel contained in a subspace]
Let $\varphi$ be a nonzero linear functional on $V$.  If a subspace
$U\subseteq V$ contains $\Ker\varphi$, then either $U=\Ker\varphi$ or
$U=V$.
\end{noteexercise}

\begin{proof}
If $x\in U\setminus\Ker\varphi$, then every $v\in V$ can be written
$v=(v-\varphi(v)x/\varphi(x))+\varphi(v)x/\varphi(x)$, with the first term
in $\Ker\varphi\subseteq U$ and the second in $U$.
\end{proof}

\begin{noteexercise}[Uniqueness in the one-dimensional extension]
In the one-dimensional Hahn--Banach step, the extension is unique exactly
when the lower and upper endpoints of the admissible interval coincide.
\end{noteexercise}

\begin{noteexercise}[Separating an open ball from the origin]
If an open ball $B(f_0,r)$ in a normed space does not contain $0$, then there
is $\varphi\in X^*$ such that
\[
  \operatorname{Re}\varphi(f)>0
  \qquad(f\in B(f_0,r)).
\]
\end{noteexercise}

\begin{proof}
Since $0\notin B(f_0,r)$, we have $\|f_0\|>r$.  By Hahn--Banach, choose
$\varphi\in X^*$ with $\|\varphi\|=1$ and
$\varphi(f_0)=\|f_0\|$; in the complex case, rotate a norming functional by a
unimodular scalar to obtain this normalization.  If $f=f_0+h$ with
$\|h\|<r$, then
\[
  \operatorname{Re}\varphi(f)
  =\|f_0\|+\operatorname{Re}\varphi(h)
  \ge \|f_0\|-|\varphi(h)|
  >\|f_0\|-r
  >0.
\]
\end{proof}
\sourcepages{68--69}

\begin{noteexercise}[Sublinear level sets and supporting functionals]
Let $p$ be sublinear on a real vector space.
\begin{enumerate}
  \item The set $\{x:p(x)\le1\}$ is convex.
  \item For every $x_0$ there is a linear functional $f$ such that
        $f\le p$ on $X$ and $f(x_0)=p(x_0)$.
\end{enumerate}
\end{noteexercise}

\begin{proof}
Convexity is immediate from positive homogeneity and subadditivity.  Define
$f_0(\alpha x_0)=\alpha p(x_0)$.  For $\alpha<0$, the inequality
$p(x_0)+p(-x_0)\ge0$ gives
$f_0(\alpha x_0)\le p(\alpha x_0)$.  Hahn--Banach extends $f_0$.
\end{proof}
\sourcepage{70}

\begin{noteexercise}[Hyperplanes]
A subspace $H$ of a vector space $X$ has codimension one if and only if
$H=\Ker f$ for some nonzero linear functional $f$.  If two nonzero
functionals have the same kernel, they are scalar multiples of each other.
\end{noteexercise}

\begin{proof}
If $X/H$ is one-dimensional, compose the quotient map with a nonzero linear
functional on $X/H$.  Conversely, the first isomorphism theorem gives
$X/\Ker f\cong\Ran f$, which is one-dimensional.  The proportionality claim
follows from the unique decomposition relative to a vector outside the common
kernel.
\end{proof}
\sourcepage{71}

\section{The dual of $C([a,b])$}

\begin{notetheorem}[Riesz--Stieltjes representation]
Every bounded linear functional $F$ on $C([a,b])$ has a representation
\[
  F(f)=\int_a^b f(t)\,\mathrm{d}\alpha(t),
\]
where $\alpha$ is a function of bounded variation on $[a,b]$.  With the
normalization $\alpha(a)=0$ and right continuity on $(a,b)$, the representing
function is unique and
\[
  \|F\|=\Var_{[a,b]}(\alpha).
\]
Equivalently, $C([a,b])^*$ is isometrically the space of finite signed or
complex regular Borel measures on $[a,b]$.
\end{notetheorem}

\begin{noteremark}
The notation $C([a,b])^*\cong BV([a,b])$ requires a normalization: adding a
constant to $\alpha$ does not change the Riemann--Stieltjes integral.  The
measure formulation avoids this ambiguity.
\end{noteremark}
\sourcepage{72}

\section{Adjoint operators and annihilators}

\begin{notedefinition}[Adjoint]
Let $T\in\Bop{X}{Y}$.  The adjoint is the bounded linear map
\[
  T^*\colon Y^*\to X^*,
  \qquad
  (T^*g)(x)=g(Tx).
\]
\end{notedefinition}

\begin{notetheorem}[Norm of the adjoint]
For every bounded operator $T$,
\[
  \|T^*\|=\|T\|.
\]
\end{notetheorem}

\begin{proof}
For $g\in Y^*$,
$\|T^*g\|\le\|g\|\|T\|$, so $\|T^*\|\le\|T\|$.  Conversely, given
$\varepsilon>0$, choose $x\in S_X$ with
$\|Tx\|>\|T\|-\varepsilon$.  Hahn--Banach gives $g\in S_{Y^*}$ with
$g(Tx)=\|Tx\|$.  Hence
$\|T^*\|\ge\|T^*g\|\ge|T^*g(x)|>\|T\|-\varepsilon$.
\end{proof}
\sourcepage{73}

\begin{notedefinition}[Annihilators]
For $M\subseteq X$, define
\[
  M^\perp=\{f\in X^*:f(x)=0\text{ for every }x\in M\}.
\]
For $N\subseteq X^*$, define
\[
  {}^\perp N=\{x\in X:f(x)=0\text{ for every }f\in N\}.
\]
\end{notedefinition}

\begin{notetheorem}[Range--kernel annihilator identities]
For $T\in\Bop{X}{Y}$,
\[
  \overline{\Ran T}={}^\perp\Ker T^*,
  \qquad
  \overline{\Ran T^*}^{\;w^*}=(\Ker T)^\perp.
\]
In particular,
\[
  (\Ran T)^\perp=\Ker T^*.
\]
\end{notetheorem}

\begin{proof}
A functional $g\in Y^*$ annihilates $\Ran T$ exactly when
$g(Tx)=0$ for every $x$, that is, exactly when $T^*g=0$.  Taking
preannihilators and using the Hahn--Banach closure theorem gives the first
identity.  The second is the dual statement with weak-$*$ closure.
\end{proof}
\sourcepage{74}

\begin{notedefinition}[Canonical embedding]
For a normed space $X$, define
\[
  J\colon X\to X^{**},
  \qquad
  (Jx)(f)=f(x).
\]
Then $J$ is a linear isometry.
\end{notedefinition}

\begin{proof}
The estimate $\|Jx\|\le\|x\|$ is immediate.  A norming functional
$f\in S_{X^*}$ with $f(x)=\|x\|$ gives the reverse inequality.
\end{proof}
\sourcepage{75}

\begin{notedefinition}[Reflexive space]
A normed space $X$ is reflexive if the canonical embedding is onto:
$J(X)=X^{**}$.
\end{notedefinition}

\begin{notetheorem}
Every reflexive normed space is Banach, and every finite-dimensional normed
space is reflexive.
\end{notetheorem}

\begin{proof}
The bidual is Banach, and $J(X)$ is closed precisely when $X$ is complete.
If $J$ is onto, $X$ is isometric to the Banach space $X^{**}$.  In finite
dimension, $\dim X=\dim X^{**}$ and the injective map $J$ is therefore onto.
\end{proof}
\sourcepage{76}

\begin{noteexercise}[Annihilators separate closed subspaces]
Distinct closed subspaces $Y_1,Y_2$ of a normed space have distinct
annihilators.  A subset $M\subseteq X$ is total, meaning
$\overline{\Span M}=X$, if and only if every $f\in X^*$ vanishing on $M$ is
zero.
\end{noteexercise}

\begin{proof}
If $x\in Y_1\setminus Y_2$, separate $x$ from the closed subspace $Y_2$.
The resulting functional belongs to $Y_2^\perp$ but not $Y_1^\perp$.  The
totality criterion is the annihilator criterion for density.
\end{proof}
\sourcepage{77}

\section{Baire category and its consequences}

\begin{notetheorem}[Baire category theorem]
Let $X$ be a complete metric space.
\begin{enumerate}
  \item $X$ is not a countable union of closed sets with empty interior.
  \item The intersection of countably many dense open subsets of $X$ is
        dense, and in particular nonempty.
\end{enumerate}
\end{notetheorem}

\begin{proof}
It is enough to prove the second statement.  Let $(G_n)$ be dense and open and
let $U$ be a nonempty open set.  Choose a closed ball
$\overline B(x_1,r_1)\subseteq U\cap G_1$.  Inductively choose
\[
  \overline B(x_{n+1},r_{n+1})
  \subseteq B(x_n,r_n)\cap G_{n+1},
  \qquad r_{n+1}<\frac12r_n.
\]
The centres form a Cauchy sequence and converge to some $x\in X$.  The nesting
puts $x$ in every $G_n$ and in $U$.  Taking complements gives the first
statement.
\end{proof}
\sourcepage{78}

\begin{notetheorem}[Cantor intersection theorem]
Let $(F_n)$ be nonempty closed subsets of a complete metric space such that
\[
  F_1\supseteq F_2\supseteq\cdots,
  \qquad
  \diam(F_n)\longrightarrow0.
\]
Then $\bigcap_nF_n$ consists of exactly one point.
\end{notetheorem}

\begin{proof}
Choose $x_n\in F_n$.  If $m,n\ge N$, then
$d(x_m,x_n)\le\diam(F_N)$, so $(x_n)$ is Cauchy.  Let $x_n\to x$.  For every
$N$, all sufficiently late terms lie in $F_N$, and closedness gives
$x\in F_N$.  Two points in the intersection have distance at most
$\diam(F_N)$ for every $N$, hence are equal.
\end{proof}

\begin{notetheorem}[A closed-set form of Baire]
If $X$ is complete and $X=\bigcup_{n=1}^{\infty}F_n$ with each $F_n$ closed,
then some $F_n$ has nonempty interior.  More generally, if an open set
$U\subseteq\bigcup_nF_n$, then some $F_n$ has nonempty interior relative to
$U$.
\end{notetheorem}
\sourcepages{79--80}

\begin{notetheorem}[No countable Hamel basis]
An infinite-dimensional Banach space has no countable Hamel basis.
Consequently, the vector space of polynomials cannot be complete under any
norm for which its usual countable monomial basis is a Hamel basis.
\end{notetheorem}

\begin{proof}
If $(e_n)$ were a Hamel basis, then
\[
  X=\bigcup_{n=1}^{\infty}\Span\{e_1,\ldots,e_n\}.
\]
Each finite-dimensional span is closed and has empty interior: a subspace
with nonempty interior is the whole space.  This contradicts Baire's theorem.
\end{proof}
\sourcepage{81}

\begin{notetheorem}[Cantor's theorem characterizes completeness]
A metric space $X$ is complete if and only if every nested sequence of
nonempty closed sets with diameters tending to zero has nonempty intersection.
\end{notetheorem}

\begin{proof}
Only the converse needs proof.  For a Cauchy sequence $(x_n)$ set
\[
  F_n=\overline{\{x_k:k\ge n\}}.
\]
The sets are nonempty, closed, nested, and their diameters tend to zero.
A point in their intersection is the limit of $(x_n)$.
\end{proof}
\sourcepage{82}

\begin{notetheorem}[Existence of nowhere differentiable functions]
There are continuous real-valued functions on $[0,1]$ that are differentiable
at no point.
\end{notetheorem}

\begin{proof}
Work in the Banach space $C([0,1])$.  For $m,n\in\mathbb N$, let $F_{m,n}$ be
the set of functions $f$ for which there is $x\in[0,1]$ satisfying
\[
  |f(y)-f(x)|\le m|y-x|
  \quad\text{whenever }0<|y-x|<1/n.
\]
Each $F_{m,n}$ is closed: for uniformly convergent $f_k\to f$, choose witness
points $x_k$, pass to a convergent subsequence, and pass the inequality to the
limit.  Each $F_{m,n}$ has empty interior.  Indeed, approximate the centre of
any ball uniformly by a polygonal function and add a sufficiently small,
high-frequency sawtooth whose slopes have absolute value larger than
$m$ plus the finitely many background slopes.  The perturbed function remains
in the ball but has no point satisfying the displayed local Lipschitz bound.

If $f$ is differentiable at $x$, then for some $m,n$ it belongs to $F_{m,n}$.
Thus the set of functions differentiable somewhere is contained in the
countable union $\bigcup_{m,n}F_{m,n}$, a meagre set.  Baire's theorem leaves
a continuous function outside this union.
\end{proof}
\sourcepages{83--84}

\begin{notetheorem}[Uniform boundedness principle]
Let $X$ be Banach, let $Y$ be normed, and let
$\mathcal A\subseteq\Bop{X}{Y}$.  If
\[
  \sup_{T\in\mathcal A}\|Tx\|<\infty
  \qquad(x\in X),
\]
then
\[
  \sup_{T\in\mathcal A}\|T\|<\infty.
\]
\end{notetheorem}

\begin{proof}
Set
\[
  F_n=\{x:\sup_{T\in\mathcal A}\|Tx\|\le n\}.
\]
The sets $F_n$ are closed and cover $X$.  Baire gives
$B(x_0,r)\subseteq F_N$ for some $x_0,r,N$.  If $\|h\|<r$, then
$x_0\pm h\in F_N$, so
\[
  2\|Th\|=\|T(x_0+h)-T(x_0-h)\|\le2N.
\]
Taking $h=(r/2)x$ for $\|x\|\le1$ gives
$\|Tx\|\le2N/r$ uniformly in $T$.
\end{proof}
\sourcepage{85}

\begin{notetheorem}[Banach--Steinhaus pointwise-limit theorem]
Let $X$ be Banach, let $Y$ be normed, and let $A_n\in\Bop{X}{Y}$.  Suppose
$A_nx$ converges for every $x\in X$, and define
$Ax=\lim_nA_nx$.  Then $A\in\Bop{X}{Y}$ and
\[
  \|A\|\le\liminf_{n\to\infty}\|A_n\|.
\]
\end{notetheorem}

\begin{proof}
For every $x$, the sequence $(A_nx)$ is bounded, so uniform boundedness gives
$M=\sup_n\|A_n\|<\infty$.  Linearity passes to the limit and
$\|Ax\|\le M\|x\|$.  For $\|x\|=1$,
\[
  \|Ax\|=\lim_n\|A_nx\|
  \le\liminf_n\|A_n\|;
\]
take the supremum over the unit sphere.
\end{proof}
\sourcepage{86}

\begin{noteexample}[The polynomial space is not complete]
Let $P([0,1])$ be the polynomial space with the supremum norm.  The Taylor
polynomials
\[
  p_n(x)=\sum_{k=0}^{n}\frac{x^k}{k!}
\]
converge uniformly on $[0,1]$ to $e^x$.  Thus $(p_n)$ is Cauchy in the
supremum norm, but its limit is not a polynomial.  Hence $P([0,1])$ is not
complete.  This direct argument also avoids the source's attempted use of
coefficient or derivative functionals, which are not bounded on the full
polynomial space with the supremum norm.
\end{noteexample}
\sourcepages{86--87}

\begin{noteexercise}[Two applications of uniform boundedness]
\begin{enumerate}
  \item If $x_n\in X$ and $f(x_n)$ is bounded for every $f\in X^*$, then
        $(x_n)$ is norm bounded.
  \item For Banach spaces $X,Y$ and $T_n\in\Bop{X}{Y}$, the following are
        equivalent:
        \begin{enumerate}
          \item $(\|T_n\|)$ is bounded;
          \item $(\|T_nx\|)$ is bounded for every $x\in X$;
          \item $(g(T_nx))$ is bounded for every $x\in X$ and $g\in Y^*$.
        \end{enumerate}
\end{enumerate}
\end{noteexercise}

\begin{proof}
For (1), apply uniform boundedness to the canonical evaluations
$Jx_n\in X^{**}$ and use $\|Jx_n\|=\|x_n\|$.  For (2), only the last
implication is nontrivial.  First apply uniform boundedness on $Y^*$ to the
functionals $g\mapsto g(T_nx)$ for each fixed $x$, obtaining pointwise norm
boundedness in $Y$; then apply it on $X$ to the operators $T_n$.
\end{proof}
\sourcepage{88}

\subsection{Quadrature, interpolation, and Fourier sums}

\begin{notetheorem}[P\'olya's theorem for quadrature formulas]
Let $w\in L^1(0,1)$ and let
\[
  L_n(f)=\sum_{j=0}^{m_n}w_j^{(n)}f(x_j^{(n)}),
  \qquad f\in C([0,1]),
\]
where, for each $n$, the nodes $x_j^{(n)}\in[0,1]$ are distinct.  Assume
\[
  L_n(p)\longrightarrow\int_0^1p(t)w(t)\dd t
  \qquad\text{for every polynomial }p.
\]
Then the same convergence holds for every $f\in C([0,1])$ if and only if
\[
  \sup_n\sum_j|w_j^{(n)}|<\infty.
\]
\end{notetheorem}

\begin{proof}
The functional norm is
$\|L_n\|=\sum_j|w_j^{(n)}|$: the upper bound is immediate, and a continuous
piecewise-linear function can be chosen with prescribed signs at the finitely
many nodes.  If the norms are uniformly bounded, approximate $f$ uniformly by
a polynomial $p$ and estimate
\[
 |L_n(f)-\textstyle\int fw|
 \le\|L_n\|\|f-p\|_\infty
     +|L_n(p)-\textstyle\int pw|
     +\|w\|_1\|p-f\|_\infty.
\]
Conversely, convergence for every $f$ makes $(L_n(f))$ pointwise bounded, so
uniform boundedness yields $\sup_n\|L_n\|<\infty$.
\end{proof}
\sourcepages{89--90}

\begin{notetheorem}[Norm of the Lagrange interpolation operator]
For distinct nodes $0\le t_0<\cdots<t_n\le1$, let
\[
  (L_nf)(t)=\sum_{k=0}^nf(t_k)p_k(t),
  \qquad
  p_k(t)=\prod_{j\ne k}\frac{t-t_j}{t_k-t_j}.
\]
Then
\[
  \|L_n\|
  =\sup_{t\in[0,1]}\sum_{k=0}^n|p_k(t)|.
\]
\end{notetheorem}

\begin{proof}
The triangle inequality gives the upper bound.  Choose $t^*$ attaining the
maximum and a continuous piecewise-linear function $f$ with
$\|f\|_\infty=1$ and
$f(t_k)=\sign p_k(t^*)$ (using arbitrary phases in the complex case).  Then
$|(L_nf)(t^*)|=\sum_k|p_k(t^*)|$.
\end{proof}
\sourcepages{91--92}

\begin{notetheorem}[Fourier partial sums and the Dirichlet kernel]
On the Banach space of continuous $2\pi$-periodic functions, the $n$th Fourier
partial-sum operator satisfies
\[
  (S_nf)(t)=\frac1{2\pi}\int_{-\pi}^{\pi}f(t-u)D_n(u)\dd u,
  \qquad
  D_n(u)=\frac{\sin((n+\tfrac12)u)}{\sin(u/2)},
\]
and
\[
  \|S_n\|=\frac1{2\pi}\int_{-\pi}^{\pi}|D_n(u)|\dd u.
\]
Moreover, $\|S_n\|\to\infty$ at logarithmic order.  Hence there exists a
continuous periodic function whose Fourier partial sums diverge at a point.
\end{notetheorem}

\begin{proof}
The convolution formula gives the upper bound.  For the reverse bound,
approximate the phase function $u\mapsto\sign(D_n(-u))$ by continuous
periodic functions, bounded in modulus by one, in the weighted space
$L^1(|D_n(u)|\dd u)$.  On each interval where the numerator has fixed sign
and $|u|$ is small, $|\sin(u/2)|\le |u|/2$.  Summing the resulting harmonic
lower bounds yields $\|S_n\|\to\infty$.  If every continuous $f$ had
bounded partial sums at zero, uniform boundedness would force
$\sup_n\|S_n\|<\infty$, a contradiction.
\end{proof}
\sourcepages{93--94}

\begin{notetheorem}[Separate continuity of bilinear maps]
Let $X$ and $Y$ be Banach spaces, let $Z$ be normed, and let
$B\colon X\times Y\to Z$ be bilinear.  If $B(\,\cdot\,,y)$ is bounded for
every fixed $y$ and $B(x,\,\cdot\,)$ is bounded for every fixed $x$, then
there is $C<\infty$ such that
\[
  \|B(x,y)\|\le C\|x\|\|y\|.
\]
Thus $B$ is jointly continuous.
\end{notetheorem}

\begin{proof}
For $y\in Y$, define $A_yx=B(x,y)$.  For each fixed $x$, the family
$\{A_y:\|y\|\le1\}$ is pointwise bounded because
$y\mapsto B(x,y)$ is bounded.  Uniform boundedness gives
$\sup_{\|y\|\le1}\|A_y\|=C<\infty$, which is the desired estimate.
\end{proof}

\begin{noteexample}[Completeness is essential]
Let $c_{00}$ carry the supremum norm and define
\[
  B(x,y)=\sum_{k=1}^{\infty}x_ky_k.
\]
The sum is finite.  For fixed $x$ or fixed $y$, the resulting functional is
bounded, but $B$ is not jointly bounded: taking $x=y=(1,\ldots,1,0,\ldots)$
with $n$ ones gives $B(x,y)=n$ while both sup norms equal one.
\end{noteexample}
\sourcepages{95--96}

\begin{notecorollary}[Positive quadrature weights]
In P\'olya's theorem, if every $w_j^{(n)}\ge0$, convergence on polynomials
already implies convergence on all continuous functions.
\end{notecorollary}

\begin{proof}
Applying polynomial convergence to the constant polynomial $1$ gives
\[
  \sum_j|w_j^{(n)}|=\sum_jw_j^{(n)}=L_n(1)
  \longrightarrow\int_0^1w(t)\dd t.
\]
The norms are therefore uniformly bounded, and P\'olya's theorem applies.
\end{proof}
\sourcepage{97}

\section{Weak and weak-star convergence}

\begin{notedefinition}[Strong and weak convergence]
A sequence $(x_n)$ in a normed space $X$ converges strongly to $x$ if
$\|x_n-x\|\to0$.  It converges weakly to $x$, written
$x_n\weakto x$, if
\[
  f(x_n)\longrightarrow f(x)
  \qquad(f\in X^*).
\]
\end{notedefinition}

\begin{notelemma}[Basic facts]
If $x_n\weakto x$, then:
\begin{enumerate}
  \item the weak limit is unique;
  \item every subsequence converges weakly to $x$;
  \item $(x_n)$ is norm bounded;
  \item $\|x\|\le\liminf_n\|x_n\|$.
\end{enumerate}
\end{notelemma}

\begin{proof}
Functionals separate points, giving uniqueness.  The subsequence statement is
immediate.  For boundedness, regard $Jx_n$ as functionals on the Banach space
$X^*$.  They are pointwise bounded, so uniform boundedness gives
$\sup_n\|Jx_n\|=\sup_n\|x_n\|<\infty$.  A norming functional for $x$ gives
\[
  \|x\|=|f(x)|=\lim_n|f(x_n)|
  \le\liminf_n\|x_n\|.
\]
\end{proof}
\sourcepages{98--99}

\begin{notetheorem}[Strong versus weak convergence]
Strong convergence implies weak convergence.  The converse holds in
finite-dimensional normed spaces, but fails in every infinite-dimensional
Hilbert space: an orthonormal sequence $(e_n)$ satisfies $e_n\weakto0$ while
$\|e_n\|=1$.
\end{notetheorem}

\begin{proof}
For strong convergence,
$|f(x_n)-f(x)|\le\|f\|\|x_n-x\|$.  In finite dimension, weak convergence of
the coordinate functionals gives convergence of all coordinates and hence
norm convergence.  In a Hilbert space, Riesz representation writes every
$f$ as $f(x)=\langle x,z\rangle$; Bessel's inequality gives
$\langle e_n,z\rangle\to0$.
\end{proof}

\begin{noteremark}[Schur's theorem]
In $\ell^1$, weak convergence of sequences implies norm convergence.  This is
a special property of $\ell^1$, not a general feature of infinite-dimensional
Banach spaces.
\end{noteremark}
\sourcepages{99--100}

\begin{notelemma}[Testing weak convergence on a total set]
Let $(x_n)$ be bounded in $X$ and let $M\subseteq X^*$ have norm-dense linear
span in $X^*$.  If $f(x_n)\to f(x)$ for every $f\in M$, then
$x_n\weakto x$.
\end{notelemma}

\begin{proof}
The convergence extends by linearity to $\Span M$.  If
$\|x_n\|,\|x\|\le C$ and $g\in X^*$, choose $f\in\Span M$ with
$\|g-f\|<\varepsilon/(3C)$.  Then
\[
 |g(x_n-x)|
 \le |(g-f)(x_n)|+|f(x_n-x)|+|(f-g)(x)|<\varepsilon
\]
for all sufficiently large $n$.
\end{proof}
\sourcepage{100}

\begin{notetheorem}[Weak convergence in $\ell^p$]
Let $1<p<\infty$.  A sequence $x^{(n)}=(x_k^{(n)})$ in $\ell^p$ converges
weakly to $x=(x_k)$ if and only if
\begin{enumerate}
  \item $(\|x^{(n)}\|_p)$ is bounded;
  \item $x_k^{(n)}\to x_k$ for every coordinate $k$.
\end{enumerate}
\end{notetheorem}

\begin{proof}
Weak convergence gives both conditions.  Conversely, coordinate evaluations
span the finitely supported sequences, which are dense in
$(\ell^p)^*=\ell^q$.  Apply the total-set criterion.
\end{proof}

\begin{notetheorem}[Radon--Riesz property in uniformly convex spaces]
If $X$ is uniformly convex, $x_n\weakto x$, and
$\|x_n\|\to\|x\|$, then $x_n\to x$ in norm.
\end{notetheorem}

\begin{proof}
The case $x=0$ is immediate.  Otherwise put
$u_n=x_n/\|x_n\|$ and $u=x/\|x\|$.  Then $u_n\weakto u$ and
\[
  \left\|\frac{u_n+u}{2}\right\|
  \ge \operatorname{Re}f\!\left(\frac{u_n+u}{2}\right)\longrightarrow1
\]
for a norming functional $f$ of $u$.  Uniform convexity forces
$\|u_n-u\|\to0$, and the convergence of norms gives $\|x_n-x\|\to0$.
\end{proof}
\sourcepage{101}

\begin{notetheorem}[Weak convergence and operators]
Let $A\in\Bop{X}{Y}$.
\begin{enumerate}
  \item If $x_n\weakto x$, then $Ax_n\weakto Ax$.
  \item If $A$ is compact and $x_n\weakto x$, then
        $\|Ax_n-Ax\|\to0$.
  \item If $B\colon X\times Y\to Z$ is bounded bilinear,
        $x_n\to x$ strongly, and $y_n\weakto y$, then
        $B(x_n,y_n)\weakto B(x,y)$; when $Z=\mathbb K$, the convergence is
        scalar.
\end{enumerate}
\end{notetheorem}

\begin{proof}
For (1), $g(Ax_n)=(A^*g)(x_n)$.  For (2), weak convergence makes $(x_n)$
bounded.  Every subsequence of $(Ax_n)$ has a norm-convergent further
subsequence, and (1) forces its limit to be $Ax$; hence the whole sequence
converges in norm.  For (3), $(y_n)$ is bounded and
\[
  \|B(x_n,y_n)-B(x,y_n)\|
  \le\|B\|\|x_n-x\|\|y_n\|\to0,
\]
while $B(x,\cdot)$ is a bounded linear map and therefore preserves weak
convergence.
\end{proof}
\sourcepage{102}

\begin{notedefinition}[Weak topology and weak-$*$ topology]
The weak topology on $X$ is the coarsest topology making every member of
$X^*$ continuous.  The weak-$*$ topology on $X^*$ is the coarsest topology
making every evaluation $f\mapsto f(x)$, $x\in X$, continuous.
\end{notedefinition}

\begin{notetheorem}[Finite dimensions and sequences]
In finite dimension, the weak and norm topologies coincide.  In an
infinite-dimensional normed space the weak topology is strictly coarser.
For sequences,
\[
  x_n\weakto x
  \quad\Longleftrightarrow\quad
  f(x_n)\to f(x)\text{ for all }f\in X^*.
\]
\end{notetheorem}
\sourcepage{103}

\begin{notedefinition}[Weak-$*$ convergence of functionals]
A sequence $(f_n)$ in $X^*$ converges weak-$*$ to $f$ if
\[
  f_n(x)\longrightarrow f(x)
  \qquad(x\in X).
\]
Weak convergence in $X^*$ implies weak-$*$ convergence.  If $X$ is reflexive,
the two notions agree on $X^*$.
\end{notedefinition}

\begin{notetheorem}[Strong operator limits]
Let $X$ be Banach, let $Y$ be normed, and let
$T_n\in\Bop{X}{Y}$.  If $T_nx$ converges in $Y$ for every $x\in X$, then the
pointwise limit $Tx=\lim_nT_nx$ belongs to $\Bop{X}{Y}$ and
$\sup_n\|T_n\|<\infty$.
\end{notetheorem}

\begin{proof}
The family $(T_n)$ is pointwise bounded, so uniform boundedness gives a common
operator-norm bound.  Linearity and boundedness pass to the pointwise limit.
\end{proof}
\sourcepage{104}

\begin{notedefinition}[Operator convergence]
For $T_n,T\in\Bop{X}{Y}$:
\begin{itemize}
  \item $T_n\to T$ uniformly if $\|T_n-T\|\to0$;
  \item $T_n\to T$ strongly if $T_nx\to Tx$ for every $x\in X$;
  \item $T_n\to T$ weakly if $T_nx\weakto Tx$ for every $x\in X$.
\end{itemize}
\end{notedefinition}

\begin{notetheorem}[Criterion for strong operator convergence]
Let $X,Y$ be Banach spaces and let $(T_n)\subseteq\Bop{X}{Y}$.  The sequence
converges strongly if and only if
\begin{enumerate}
  \item $\sup_n\|T_n\|<\infty$;
  \item for every $x$ in a total subset $M\subseteq X$, the sequence
        $(T_nx)$ is Cauchy in $Y$.
\end{enumerate}
\end{notetheorem}

\begin{proof}
Only sufficiency needs proof.  Cauchy convergence extends by linearity to
$\Span M$.  Given $x\in X$, approximate it by $y\in\Span M$ and use the
uniform bound to control
$\|(T_n-T_m)(x-y)\|$, while $(T_ny)$ is Cauchy.  Completeness of $Y$ supplies
the pointwise limit.
\end{proof}

\begin{notecorollary}
A sequence $(f_n)\subseteq X^*$ converges weak-$*$ if and only if it is norm
bounded and $f_n(x)$ is Cauchy for every $x$ in a total subset of $X$.
\end{notecorollary}
\sourcepage{105}

\begin{notetheorem}[Reflexivity results recorded in the source]
The following standard results are used later.
\begin{enumerate}
  \item Every finite-dimensional normed space is reflexive.
  \item Every Hilbert space is reflexive.
  \item Closed subspaces and finite products of reflexive Banach spaces are
        reflexive.
  \item $L^p(\Omega)$ is reflexive for $1<p<\infty$.
  \item Every uniformly convex Banach space is reflexive
        (Milman--Pettis theorem).
  \item A Banach space is reflexive if and only if every bounded sequence has
        a weakly convergent subsequence (Eberlein--\v{S}mulian plus weak
        compactness of the unit ball).
\end{enumerate}
\end{notetheorem}

\begin{proof}[Editorial note on depth]
Milman--Pettis and Eberlein--\v{S}mulian are deep theorems.  The source records
them without complete proofs; they are therefore stated accurately here and
not represented as consequences of the short handwritten sketches.
\end{proof}
\sourcepage{106}

\begin{notetheorem}[Completeness and the canonical image]
For a normed space $X$, the canonical image $J(X)$ is closed in $X^{**}$ if
and only if $X$ is Banach.
\end{notetheorem}

\begin{proof}
If $X$ is Banach, its isometric image is complete and hence closed.  Conversely,
if $J(X)$ is closed, it is a Banach subspace of $X^{**}$, and the isometry
$J\colon X\to J(X)$ makes $X$ complete.
\end{proof}
\sourcepage{107}

\begin{notelemma}[Dense testing in the dual]
Let $(x_n)$ be bounded in a normed space $X$ and let $Y\subseteq X^*$ be
norm dense.  If $y(x_n)\to y(x)$ for every $y\in Y$, then $x_n\weakto x$.
\end{notelemma}

\begin{notetheorem}[Bounded almost-everywhere convergence in $L^p$]
Let $\Omega\subseteq\mathbb R^d$ be open, let $1<p<\infty$, and let
$(f_n)$ be bounded in $L^p(\Omega)$.  If $f_n\to f$ almost everywhere and
$f\in L^p(\Omega)$, then $f_n\weakto f$ in $L^p(\Omega)$.
\end{notetheorem}

\begin{proof}
Let $q$ be conjugate to $p$.  Given $g\in L^q(\Omega)$, first approximate it
in $L^q$ by a bounded function $g_0$ supported on a set of finite measure.
The common $L^p$ bound controls the error uniformly in $n$.  On the support of
$g_0$, almost-everywhere convergence together with the common $L^p$ bound
implies $f_n-f\to0$ in $L^1$: use Egorov's theorem off a set of arbitrarily
small measure and H\"older's inequality on that exceptional set.  Therefore
$\int g_0(f_n-f)\to0$, and the approximation gives
$\int g(f_n-f)\to0$.
\end{proof}

\begin{noteexample}[The alternating-cell sequence in $L^2(0,1)$]
For $n\ge1$ and $0\le j\le n-1$, define
\[
 f_n(t)=
 \begin{cases}
   0,&\displaystyle \frac jn\le t<\frac jn+\frac1{2n},\\[4pt]
   1,&\displaystyle \frac jn+\frac1{2n}\le t<\frac{j+1}{n}.
 \end{cases}
\]
Then $f_n\weakto\tfrac12$ in $L^2(0,1)$.  Indeed, the integrals of
$f_n-\tfrac12$ over every interval $[0,t]$ tend to zero, and finite linear
combinations of interval indicators are dense in $L^2(0,1)$; the sequence is
bounded, so the dense-testing lemma applies.  The convergence is not strong,
because
\[
  \left\|f_n-\frac12\right\|_2=\frac12
  \qquad(n\ge1).
\]
\end{noteexample}
\sourcepages{108--109}

\begin{notetheorem}[Weak lower semicontinuity of the norm]
If $x_n\weakto x$, then
\[
  \|x\|\le\liminf_{n\to\infty}\|x_n\|.
\]
\end{notetheorem}

\begin{proof}
Choose $f\in X^*$ with $\|f\|=1$ and $f(x)=\|x\|$.  Then
$|f(x_n)|\le\|x_n\|$ and $f(x_n)\to f(x)$, which gives the inequality.
\end{proof}

\begin{noteexercise}[Weak limits lie in closed spans]
If $x_n\in Y$ for a linear subspace $Y\subseteq X$ and $x_n\weakto x$, then
$x\in\overline Y$.
\end{noteexercise}

\begin{proof}
If $x\notin\overline Y$, Hahn--Banach gives $f\in X^*$ vanishing on $Y$ but
not at $x$, contradicting $f(x_n)\to f(x)$.
\end{proof}
\sourcepage{109}

\begin{notetheorem}[Weak-$*$ sequential compactness in the separable case]
Let $X$ be separable and let $(f_n)$ be a bounded sequence in $X^*$.  Then
$(f_n)$ has a weak-$*$ convergent subsequence.
\end{notetheorem}

\begin{proof}
Choose a dense sequence $(x_j)$ in $X$.  Repeatedly extract subsequences so
that $f_n(x_j)$ converges for each $j$, then take the diagonal subsequence.
The pointwise limit on the dense span is bounded by the common norm bound and
therefore extends uniquely to a member of $X^*$.  Approximation from the dense
span shows weak-$*$ convergence on all of $X$.
\end{proof}

\begin{noteremark}
This is the sequential form of Banach--Alaoglu for separable preduals.  In
nonseparable settings, weak-$*$ compactness need not be sequential compactness.
\end{noteremark}
\sourcepage{110}

\begin{notetheorem}[Toeplitz limit theorem]
Let $A=(a_{nk})_{n,k\ge1}$ be an infinite matrix and define
\[
  (Ax)_n=\sum_{k=1}^{\infty}a_{nk}x_k
\]
whenever the sum is meaningful.  The method is regular, meaning it preserves
limits of convergent scalar sequences, if and only if
\[
  \sup_n\sum_{k=1}^{\infty}|a_{nk}|<\infty,
  \qquad
  a_{nk}\longrightarrow0\quad(n\to\infty)\text{ for each fixed }k,
\]
and
\[
  \sum_{k=1}^{\infty}a_{nk}\longrightarrow1.
\]
\end{notetheorem}

\begin{proof}[Proof sketch]
The necessity follows by testing the method on the constant sequence and on
the coordinate sequences, while uniform boundedness on the Banach space $c$
of convergent sequences gives the uniform row-$\ell^1$ bound.  Conversely,
write $x_k=L+(x_k-L)$, split the remainder into a finite head and a uniformly
small tail, and use the three conditions.
\end{proof}
\sourcepage{111}

\section{Open mapping and bounded inverse theorems}

\begin{notetheorem}[Open mapping theorem]
Let $X$ and $Y$ be Banach spaces and let $T\in\Bop{X}{Y}$ be surjective.
Then $T$ is an open map.
\end{notetheorem}

\begin{proof}
It is enough to show that $T(B_X)$ contains a neighbourhood of $0$.  Since
\[
  Y=\bigcup_{n=1}^{\infty}T(nB_X),
\]
Baire's theorem gives $n$ such that $\overline{T(nB_X)}$ has nonempty
interior.  Taking differences and rescaling yields a number $\delta>0$ with
\[
  B_Y(0,\delta)\subseteq\overline{T(B_X)}.
\]
After rescaling, assume $B_Y(0,1)\subseteq\overline{T(B_X)}$.  Given
$y\in B_Y(0,1/2)$, choose inductively $x_k\in2^{-k}B_X$ so that
\[
  \left\|y-T\sum_{j=1}^{k}x_j\right\|<2^{-k-1}.
\]
The series $\sum_kx_k$ converges to $x\in B_X$, and continuity gives
$Tx=y$.  Thus a ball around $0$ lies in $T(B_X)$; translation and scaling
prove openness.
\end{proof}
\sourcepages{112--113}

\begin{notecorollary}[Bounded inverse theorem]
A bounded bijective linear map between Banach spaces has a bounded inverse.
\end{notecorollary}

\begin{proof}
The inverse is continuous because the original map is open.
\end{proof}

\begin{noteexample}[A priori estimate for a boundary-value problem]
Let
\[
  L\colon C^2_0([0,1])\to C([0,1]),
  \qquad
  Lu=a u''+b u'+cu,
\]
where $C^2_0([0,1])=\{u\in C^2([0,1]):u(0)=u(1)=0\}$ carries the norm
$\|u\|_\infty+\|u'\|_\infty+\|u''\|_\infty$.  If the two-point
boundary-value problem $Lu=f$ has a unique solution for every
$f\in C([0,1])$, then $L$ is a bounded bijection between Banach spaces.
Therefore there is $C>0$ such that
\[
  \|u\|_\infty+\|u'\|_\infty+\|u''\|_\infty
  \le C\|Lu\|_\infty.
\]
\end{noteexample}
\sourcepage{113}

\begin{notetheorem}[Equivalence of complete norms]
Let $\|\cdot\|_1$ and $\|\cdot\|_2$ be two norms on the same vector space
$X$, and assume both make $X$ Banach.  If
\[
  \|x\|_2\le C\|x\|_1
\]
for all $x$, then the norms are equivalent.
\end{notetheorem}

\begin{proof}
The identity map from $(X,\|\cdot\|_1)$ to $(X,\|\cdot\|_2)$ is a bounded
bijection.  The bounded inverse theorem gives
$\|x\|_1\le C'\|x\|_2$.
\end{proof}
\sourcepage{114}

\begin{notetheorem}[Stability of surjectivity]
Let $X,Y$ be Banach spaces and let $A\in\Bop{X}{Y}$ be surjective.  Then
there is $\varepsilon>0$ such that every
$\widetilde A\in\Bop{X}{Y}$ with
$\|\widetilde A-A\|<\varepsilon$ is surjective.
\end{notetheorem}

\begin{proof}
The open mapping theorem gives $c>0$ such that every $y\in Y$ has a preimage
$x$ with $Ax=y$ and $\|x\|\le c\|y\|$.  Choose
$\varepsilon<1/c$.  For a target $y$, construct $x_0,x_1,\ldots$ with
\[
  Ax_0=y,
  \qquad
  Ax_{n+1}=(A-\widetilde A)x_n,
  \qquad
  \|x_{n+1}\|\le c\|A-\widetilde A\|\|x_n\|.
\]
The series $x=\sum_nx_n$ converges and telescoping gives
$\widetilde Ax=y$.
\end{proof}
\sourcepage{115}

\begin{notetheorem}[Closed range and bounded inverse on the range]
Let $X,Y$ be Banach spaces and let $T\in\Bop{X}{Y}$ be injective.  The
inverse $T^{-1}\colon\Ran T\to X$ is bounded if and only if $\Ran T$ is
closed in $Y$.
\end{notetheorem}

\begin{proof}
If $T^{-1}$ is bounded and $Tx_n\to y$, then $(x_n)$ is Cauchy and converges
to $x$; hence $y=Tx\in\Ran T$.  Conversely, if $\Ran T$ is closed, it is
Banach, and the bounded bijection $T\colon X\to\Ran T$ has bounded inverse.
\end{proof}
\sourcepage{116}

\section{Closed graph theorem and closed operators}

\begin{notetheorem}[Products of Banach spaces]
If $X$ and $Y$ are Banach spaces, then $X\times Y$ is Banach under either of
the equivalent norms
\[
  \|(x,y)\|_1=\|x\|+\|y\|,
  \qquad
  \|(x,y)\|_\infty=\max\{\|x\|,\|y\|\}.
\]
\end{notetheorem}
\sourcepage{117}

\begin{notedefinition}[Closed linear operator]
Let $T\colon\Dom(T)\subseteq X\to Y$ be linear.  Its graph is
\[
  \Graph(T)=\{(x,Tx):x\in\Dom(T)\}\subseteq X\times Y.
\]
The operator is \emph{closed} if its graph is closed.
\end{notedefinition}

\begin{notetheorem}[Closed graph theorem]
Let $X,Y$ be Banach spaces.  An everywhere-defined linear map
$T\colon X\to Y$ is bounded if and only if its graph is closed.
\end{notetheorem}

\begin{proof}
Boundedness plainly implies closedness.  Conversely, if the graph is closed,
it is a Banach space with the graph norm
\[
  \|x\|_T=\|x\|+\|Tx\|.
\]
The projection $S\colon\Graph(T)\to X$, $S(x,Tx)=x$, is a bounded bijection.
The bounded inverse theorem makes $S^{-1}$ bounded, so
\[
  \|x\|+\|Tx\|=\|S^{-1}x\|_T\le C\|x\|.
\]
Thus $T$ is bounded.
\end{proof}
\sourcepages{117--119}

\begin{notetheorem}[Sequential criterion]
A linear operator $T\colon\Dom(T)\subseteq X\to Y$ is closed if and only if
whenever
\[
  x_n\in\Dom(T),
  \qquad
  x_n\to x,
  \qquad
  Tx_n\to y,
\]
one has $x\in\Dom(T)$ and $Tx=y$.
\end{notetheorem}

\begin{notetheorem}[Closedness from a closed domain]
If $\Dom(T)$ is closed in $X$ and $T\colon\Dom(T)\to Y$ is bounded, then
$T$ is closed.  Conversely, if $X$ and $Y$ are Banach and $T$ is closed,
then $\Dom(T)$ is complete for the graph norm, not necessarily for the
original norm.
\end{notetheorem}

\begin{noteexample}[Differentiation]
Let
\[
  D\colon C^1([0,1])\subset C([0,1])\to C([0,1]),
  \qquad Du=u'.
\]
With the supremum norm on both copies of $C([0,1])$, the operator is closed
but unbounded.  Indeed, $u_n\to u$ uniformly and $u_n'\to v$ uniformly imply
\[
  u_n(t)-u_n(0)=\int_0^tu_n'(s)\dd s
  \longrightarrow\int_0^tv(s)\dd s,
\]
so $u\in C^1$ and $u'=v$.  On the other hand,
$u(t)=t^n$ has norm one while $\|u'\|_\infty=n$.
\end{noteexample}
\sourcepage{120}

\begin{noteexercise}[Inverse and compact-set properties]
Let $T\colon\Dom(T)\subseteq X\to Y$ be closed.
\begin{enumerate}
  \item If $T^{-1}$ exists, then $T^{-1}$ is closed.
  \item The image $T(C)$ of a compact set $C\subseteq\Dom(T)$ is closed in
        $Y$.
  \item The inverse image $T^{-1}(K)$ of a compact set $K\subseteq Y$ is
        closed in $X$.
\end{enumerate}
\end{noteexercise}

\begin{proof}
The graph of $T^{-1}$ is obtained from the graph of $T$ by swapping the two
coordinates.  For (2), a convergent sequence $Tx_n\to y$ with $x_n\in C$
has a convergent subsequence $x_{n_k}\to x\in C$; closedness gives $Tx=y$.
For (3), if $x_n\to x$ and $Tx_n\in K$, compactness gives a subsequence
$Tx_{n_k}\to y\in K$, and closedness gives $Tx=y$.
\end{proof}
\sourcepage{121}

\begin{notetheorem}[Compact domain]
Let $K$ be a compact metric space, let $Y$ be a normed space, and let
$T\colon K\to Y$ be a bijection whose graph is closed in $K\times Y$.
Then $T^{-1}\colon T(K)\to K$ is continuous.
\end{notetheorem}

\begin{proof}
If $y_n\to y$ in $T(K)$ and $x_n=T^{-1}y_n$, compactness gives a
convergent subsequence $x_{n_k}\to x\in K$.  Closedness gives $Tx=y$, so
$x=T^{-1}y$.  Every subsequence has a further subsequence converging to the
same point $T^{-1}y$, hence the whole sequence converges.
\end{proof}

\begin{notetheorem}[Closed range from a bounded inverse]
Let $T$ be closed, let $X$ be Banach, and suppose
$T^{-1}\colon\Ran T\to X$ exists and is bounded.  Then $\Ran T$ is closed.
\end{notetheorem}

\begin{proof}
If $y_n=Tx_n\to y$, boundedness of $T^{-1}$ makes $(x_n)$ Cauchy, hence
$x_n\to x\in X$.  Closedness gives $x\in\Dom(T)$ and $Tx=y$.
\end{proof}
\sourcepage{122}

\begin{notetheorem}[Termwise differentiation]
Suppose $u_n\in C^1([0,1])$, the series $\sum_nu_n$ converges uniformly to
$u$, and the derivative series $\sum_nu_n'$ converges uniformly to $v$.
Then $u\in C^1([0,1])$ and $u'=v$.
\end{notetheorem}

\begin{proof}
The partial sums $s_N$ satisfy
$(s_N,Ds_N)\to(u,v)$ in $C([0,1])\times C([0,1])$.  Closedness of the
differentiation operator gives the conclusion.
\end{proof}

\begin{notetheorem}[Closed graph extension criterion]
Let $G\subseteq X\times Y$ be a linear subspace.  There is a linear operator
$T$ with $G=\Graph(T)$ if and only if
\[
  (0,y)\in G\quad\Longrightarrow\quad y=0.
\]
If $G$ is closed, the resulting operator is closed.
\end{notetheorem}

\begin{proof}
The condition ensures that for each $x$ there is at most one $y$ with
$(x,y)\in G$.  Define the domain as the set of such $x$ and set $Tx=y$.
Linearity follows because $G$ is a subspace, and closedness of $G$ is exactly
closedness of the operator.
\end{proof}
\sourcepages{123--124}
